\section{Scope, Method, and Qualifications}

\subsection*{Series Placement and Governing Question}

This essay is Part~I of \emph{The Creature Before God}, examining what happens when a rational creature reflexively apprehends that its act of being is continuously received from an uncreated source receiving nothing in return, why that arrest need not harden into recoil, and how consent reaches the natural-theological judgment that God loves all existing things. The duties of worship and gratitude opened at the threshold of adoration belong to Part~II (\emph{The Due Return}); supernatural revelation, Trinitarian communion, and grace belong to Part~III.

\subsection*{Epistemic Order, School Boundaries, and Exclusions}

Sections~1--3 argue on philosophical grounds alone without importing revealed premises; Section~4 extends the metaphysics to conservation and immanence before turning explicitly to Scripture, the Fathers, and the councils; Sections~5--8 integrate Augustinian phenomenology, Thomistic metaphysics of relation and virtue, apophatic restraint, and conciliar guardrails before halting at natural reason's boundary (\emph{ST} I, q.~32, a.~1; \emph{Dei Filius}, ch.~2, can.~1 \emph{De revelatione}, and ch.~4).

\begin{itemize}
\item \textbf{Editorial synthesis vocabulary:} The terms \emph{ontological vertigo}, \emph{arrest}, \emph{recoil}, and \emph{consent} are editorial synthesis grounded in Augustine (\emph{Confessions} VII.10--17), Gregory of Nazianzus (\emph{Or.}~28.21; 38.7--8), Irenaeus (\emph{Adv.\ Haer.}~IV.39.2), and Aquinas (\emph{ST} I, q.~13, a.~7; q.~105, a.~5; II-II, qq.~129, 161), not magisterial categories.
\item \textbf{Thomistic thesis vs.\ dogma:} The real distinction of essence and \emph{esse} is adopted as a Thomistic philosophical framework rather than binding dogma, noting Henry of Ghent, Scotist, and Su\'{a}rezian alternatives.
\item \textbf{Deliberate exclusions:} Following Aquinas (\emph{ST} I, q.~46, a.~2; q.~32, a.~1), \emph{Kalam} temporal-origin proofs, Big Bang inferences, and natural demonstrations of the Trinity are excluded.
\end{itemize}

\subsection*{Editions, Translations, and Rights Boundaries}

Every verbatim quotation is verified directly against public-domain editions:

\begin{itemize}
\item \textbf{Sacred Scripture:} Challoner Douay-Rheims (English) and Clementine Vulgate (Latin), keeping Vulgate/Douay-Rheims numbering (Wis.~11:25--27; Ps.~99:3, 103:29--30).
\item \textbf{St.\ Thomas Aquinas:} English Dominican \emph{Summa Theologiae} (Benziger / Project Gutenberg); \emph{De Ente et Essentia} (ch.~4--5) checked against Baur--Koch Latin and rendered in unquoted prose.
\item \textbf{Church Fathers:} Unrevised historical printings rather than modernized web revisions: Irenaeus (\emph{Adv.\ Haer.}, ANF~1), Athanasius (\emph{Contra Gentes}, NPNF2~4), Augustine (\emph{Confessions}, Pilkington, NPNF1~1; \emph{City of God}, Dods, 1871; \emph{De Genesi ad litteram}, Migne PL~34 Latin rendered in unquoted prose), Gregory of Nazianzus (\emph{Orations} 28, 38, NPNF2~7), John of Damascus (\emph{De Fide Orthodoxa} I.9, Salmond, NPNF2~9), and Pseudo-Dionysius (\emph{Celestial Hierarchy}, Parker, 1899).
\item \textbf{Councils and Catechesis:} Fourth Lateran (\emph{Firmiter credimus}, DS~800, historical Latin rendered in unquoted prose); Vatican~I (\emph{Dei Filius}, ch.~1, ch.~2, ch.~4, can.~5 \emph{De Deo rerum omnium creatore}, can.~1 \emph{De revelatione}, Latin and Manning 1912 English). The \emph{Catechism of the Catholic Church} (nn.~300--301) is paraphrased doctrinally without verbatim quotation.
\item \textbf{Other philosophical authors:} Plato, Aristotle (except inside Aquinas's \emph{Summa}), Henry of Ghent, Duns Scotus, Su\'{a}rez, Hume, and Russell are analyzed in original prose without verbatim quotations.
\end{itemize}

\subsection*{Study-Aid Notice}

This independent study aid carries no ecclesiastical imprimatur or magisterial authority, is not a critical edition of any ancient or medieval text, and reflects internal source and build verification only.
